\section{结论与未来工作}

本文研究了 Attention Residuals 从大语言模型到多智能体强化学习的迁移适配性。
实验结果表明，重型 Full/Block AttnRes 直接插入浅层 recurrent MARL agent 后并不能稳定提升性能，并且显著增加训练时间。
轻量化 AttnRes-L2 在部分指标上显示出 AUC 和 best-win 信号，完整 5M \code{5m\_vs\_6m} 对照中也获得了小幅 final-win 均值提升，但 AUC 增益很小且 seed 间方差较大。
因此，direct depth-wise transfer 并非完全无效，但其收益弱、不稳定，并不足以支撑稳定性能改进的结论。
进一步的 AttnComm-QMIX-L2 初步扩展实验显示，将 selective attention 放到 agent-wise communication 维度后，\code{5m\_vs\_6m} 上相对 QMIX baseline 出现了更明确的 final win 和 AUC 正向信号，其中 other-only communication 的平均提升更明显。

因此，当前结果不支持“Attention Residuals 可以直接作为浅层 MARL agent 后处理模块稳定提升性能”的强结论。
更合理的适配路径是将 selective attention 思想迁移到 agent-wise communication 中，用于建模智能体间协同关系。
从实验定位看，本文的主要价值在于给出一组可复现实验和适配性分析，而不是提出一个追求最优性能的新 MARL 算法。

当前实验已经足以支持上述谨慎结论，因此本文不将继续补跑 \code{3s5z}、扩展 AttnComm 到 IQL/VDN 或新增 hard map 作为定稿前的必要条件。
后续工作可以从以下方向增强证据：

\begin{itemize}
  \item 在更多 SMAC 地图上验证不同复杂度任务中的迁移表现，优先选择比 \code{3s5z} 更有区分度的中高难地图。
  \item 在 \code{5m\_vs\_6m} 上增加 paired seeds，进一步检验 AttnComm 与 AttnRes-L2 的 seed 稳定性。
  \item 在更有区分度的 SMAC 地图上继续验证 AttnComm 信号，而不仅依赖接近饱和的 \code{3s5z} sanity check。
  \item 在更多 SMAC 地图上比较 \code{qmix\_attncomm\_l2\_other} 与 \code{qmix\_attncomm\_l2\_self}，检验 other-only communication 信号是否稳定。
  \item 在 IQL 或 VDN 等其他后端上验证 AttnComm 是否仍具有类似趋势。
  \item 分析 communication attention weights 与 agent coordination 的关系。
  \item 增加统计显著性分析。
\end{itemize}
